\scanpage{34}\mspage{3}

Notons la formule fondamentale

\[
\boxed{\sigma_0(\beta)(e_0(\beta))=\deleted{\unsure{\bigl(\begin{smallmatrix}1&0\\-\beta^{-1}&1\end{smallmatrix}\bigr)}=}\deleted{\unsure{u}}e_1(-\beta^{-1})}\tag*{(25)}
\]

% slanted
\margin{d'où plus généralement $\sigma_0(\beta)(e_0(\lambda\beta))=e_1(-\lambda\beta^{-1})$}

qui \deleted{\unsure{qui}} montre en particulier que

\[
\begin{aligned}
 \deleted{\unsure{tout|pour}}\forall\sigma\in N^-,\quad &\text{on a }\sigma(U_0)=\deleted{\unsure{u}}U_1\\
 &\text{d'où }\sigma(U_1)=U_0\qquad\deleted{\text{\unsure{pa}}}
\end{aligned}\tag*{(26)}
\]

(puisque $\sigma^2=\unsure{\omega|w}$ est central).

Notons que

\[
g=\bigl(\begin{smallmatrix}a&b\\c&d\end{smallmatrix}\bigr)\in H\quad\Longrightarrow\quad g^{-1}=\bigl(\begin{smallmatrix}d&-b\\-c&a\end{smallmatrix}\bigr)\tag*{(27)}
\]

et

\[
g e_0(\beta)g^{-1}=\bigl(\begin{smallmatrix}1-ac\beta&a^2\beta\\-c^2\beta&1+ac\beta\end{smallmatrix}\bigr)\tag*{(28)}
\]

donc on a $g(U_0)\subset U_1$ sss on a

\[
\begin{aligned}
 &a^2\beta=0,\ ac\beta=0\qquad\forall\,\beta\in L_0\\
 \text{i.e.}\qquad &a^2=0,\ ac=0\\
 \text{d'où}\qquad &a=a(\underbrace{ad-bc}_{1})=0
\end{aligned}
\]

donc

\[
\left\{\begin{aligned}
 \operatorname{Transp}_H(U_0,U_1)&=\left\{\bigl(\begin{smallmatrix}a&b\\c&d\end{smallmatrix}\bigr)\in H\mid a=0\right\},\ \text{de même}\\
 \operatorname{Transp}_H(U_1,U_0)&=\left\{\bigl(\begin{smallmatrix}a&b\\c&d\end{smallmatrix}\bigr)\in H\mid d=0\right\}.
\end{aligned}\right.\tag*{(29)}
\]

% slanted
\margin{il suffit même qu'un $g\in H$ soit tel que pour \uline{un} $\beta\in L_0^*$, on ait $g(\beta)\in U_1$, pour avoir $g\in\operatorname{Transp}_H(U_0\unsure{,|\to}U_1)$}

\deleted{\unsure{et}} \deleted{Je dis que} En fait, \unsure{le|les} \deleted{\unsure{on voit}} \inserted{calcul} \deleted{\unsure{on}} montre que dans ces cas là, le Transporteur large est strict, car si $a=0$, la formule (28) donne

\[
g e_0(\beta)g^{-1}=\bigl(\begin{smallmatrix}1&0\\-c^2\beta&1\end{smallmatrix}\bigr)\tag*{$(28')$}
\]

or $ad-bc=-bc=1$, donc $b$ et $c$ sont inversibles, donc $c^2$ inversible, donc $\operatorname{int}(g)$ induit un \uline{iso} $U_0\simeq U_1$. Des formules (29) on tire

% archive sheet number: 33
